Um grafo planar com $k$ componentes conexas satisfaz
$$V-A+F = 1+k$$
Se $V \geq 3$, o número máximo de arestas de um grafo planar é $3V-6$
Dá pra conseguir isso em qualquer grafo planar em que as faces são
limitadas por triângulos. Ou seja, $A=O(V)$ em grafos planares. O número
máximo de faces é $2V-4$. Todo grafo planar tem um vértice de grau $5$
ou menos.

\subsubsection{Encontrar faces}

$O(n \log n)$

\inccode{geo/find-faces.cpp}


\subsubsection{Construir grafo planar a partir de segmentos}
$O(n^2 \log n)$

\inccode{geo/planar-graph-from-segments.cpp}

\subsubsection{Point location}

Dada uma subdivisão do plano, determinar que face da subdivisão
o ponto pertence em $O(\log n)$.

Implementado pra int. Mas dá pra mudar pra double mudando as funções
de comparação. Assume que a subdivisão está armazenada corretamente numa
DCEL e a face externa está numerada com $-1$. 

Para cada query, é retornado um pair $(1, i)$ se o ponto está na face de
número $i$ e um pair $(0, i)$ é retornado caso esteja em cima de uma edge
de número $i$.

\inccode{geo/point-location.cpp}
